% The degenerate region, which is where a reader expects the recursion to break:
% a wall across a one-cell-wide region has a single crossing, and the door is
% forced onto it, so the deletion loop runs once and does nothing.
\begin{figure}[htbp]
  \centering
  \begin{tikzpicture}[scale=1.0]
    \foreach \y in {0,1,2,3} \node[node vertex] (c\y) at (0,-\y) {};
    \foreach \y/\yn in {0/1,1/2,2/3} \draw[door] (c\y) -- (c\yn);
    % the wall the call places, at row boundary j = b_0+1
    \draw[wall] (-0.75,-1.5) -- (-0.18,-1.5);
    \draw[wall] (0.18,-1.5)  -- (0.75,-1.5);
    \node[font=\tiny, wallred, left]  at (-0.8,-1.5) {wall at $j$};
    \node[font=\tiny, doorgreen, right] at (0.8,-1.5)
         {the only crossing, and the door};
    \node[font=\tiny, black!70, right] at (0.35,-0.3)
         {$d_x=0<d_y$, so $t=0$};
    \node[font=\tiny, black!70, right] at (0.35,-2.7)
         {$i\leftarrow\textsc{RandomIn}(a_0,a_0)=a_0$: forced};
    \node[font=\tiny, black!70, below] at (0,-3.55)
         {loop runs once, with $k=i$: nothing is deleted};
  \end{tikzpicture}
  \caption{The degenerate region, where a recursion of this shape is usually
  expected to fail. A region one cell wide is cut across its other dimension,
  so the wall exists and the two halves are non-empty; but the wall has a
  single crossing edge, and the door is forced onto it because $i$ is drawn
  from a range of one value. The loop body therefore never runs, the corridor
  descends to its leaves with every door intact, and no call can disconnect it.
  The same argument with the roles of $x$ and $y$ exchanged covers a region one
  cell high.}
  \label{fig:corridor-base-case}
\end{figure}
